Come discusso in precedenza l'utilizzo concorrenziale di risorse condivise da parte di due o più Processi può portare ad inconsistenze, per cui è buona norma utilizzare oltre alle Tecniche di Controllo del Flusso di Esecuzione illustrate poc'anzi, anche Tecniche di Sincronizzazione tra i Processi in questione in maniera tale da garantire una corretta esecuzione del Programma ed una corretta gestione delle suddette Risorse comuni. \par
\section{Il Costrutto \textit{Synchronized}}
Java mette a disposizione del Programmatore un costrutto \textit{ad hoc} chiamato ``\textit{Synchronized}'', il quale va a simulare l'utilizzo di un Monitor simile a quello definito in Concurrent Pascal illustrato in precedenza per la risorsa che viene protetta. Quest'ultima può essere un metodo o un blocco di codice, indicando quindi una Sezione Critica.
\begin{figure}[H]
    \centering
    \begin{subfigure}[t]{.45\textwidth}
        \begin{lstlisting}[language=Java]
public synchronized void methodName() {
    // Metodo Protetto
}
        \end{lstlisting}
        \caption{Esempio di dichiarazione di un Metodo \textit{Synchronized}.}
        \label{subfig:SyncronizedMethod}
    \end{subfigure}
    \hfill
    \begin{subfigure}[t]{.45\textwidth}
        \begin{lstlisting}[language=Java]
synchronized(Object) {
    // Blocco di Codice Protetto
}
        \end{lstlisting}
        \caption{Esempio di dichiarazione di un Blocco di Codice \textit{Synchronized}.}
        \label{subfig:SyncronizedBlock}
    \end{subfigure}
    \caption[Esempio di utilizzo del Costrutto \textit{Synchronized}]{Esempio di utilizzo del Costrutto \textit{Synchronized}.}
    \label{fig:SyncronizedBasicExample}
\end{figure}
Nel caso della \autoref{subfig:SyncronizedMethod} si parla di Sincronizzazione Implicita, ed un solo \textit{Thread} alla volta potrà eseguire il metodo protetto. Nel caso della \autoref{subfig:SyncronizedBlock} si parla invece di Sincronizzazione Esplicita ed ancora una volta in ogni momento un solo Processo Leggero potrà eseguire tale blocco protetto. \par
All'interno di un Metodo \textit{Synchronized} è poi possibile utilizzare due speciali primitive:
\begin{itemize}
    \item \textbf{\textit{wait()}} \par
    Il funzionamento è analogo all'omonima primitiva per un Monitor in Concurrent Pascal: permette infatti di sospendere il Processo chiamante su una coda associata all'oggetto dotato del suddetto metodo \textit{Synchronized}, rilasciando quindi il Mutex ad esso connesso. Essendo analogo anche nel funzionamento al costrutto Monitor del Concurrent Pascal, una chiamata del Metodo \texttt{wait()} in un momento in cui nessun altro processo è sospeso viene persa, non avendo alcun effetto. 
    \item \textbf{\textit{notify()}} \par
    Questo ha invece un comportamento assimilabile a quello della primitiva \texttt{signal()} sempre disponibile nel contesto di un Monitor in Concurrent Pascal e permette quindi di risvegliare il primo dei \textit{Thread} sospesi sulla coda dell'Oggetto sul quale è stato invocato il metodo. \par
    Ne esiste una variante ``\texttt{notifyAll()}'' che risveglia tutti i \textit{Thread} ``addormentati'' sulla coda dell'Oggetto. Si noti che comunque un solo processo alla volta potrà riacquisire il \textit{Lock} vista la sua natura mutualmente esclusiva.
\end{itemize}
\section{Il costrutto \textit{Lock}}